\documentclass[11pt]{article}
\usepackage[margin=1in]{geometry}
\usepackage{amsmath,amssymb}
\usepackage{bm}
\usepackage{booktabs}
\usepackage{xcolor}

% --- Notation macros (consistent with meta_simulator.tex) ---
\newcommand{\rmB}{\mathbf{B}}
\newcommand{\rmW}{\mathbf{W}}
\newcommand{\clF}{\mathcal{F}}
\newcommand{\clM}{\mathcal{M}}
\providecommand{\cref}[1]{Eq.~(\ref{#1})}

\title{Extending a Trained Model-Class CogFormer to a New Family\\
       via Fine-Tuning: the LBA Experiment}
\author{CogFormer-2}
\date{}

\begin{document}
\maketitle

%============================================================================
\section{Motivation}
%============================================================================

The base model-class CogFormer is trained as a single amortized posterior
estimator over a \emph{model class} $\clM$ comprising three sequential-sampling
families: the drift-diffusion model (DDM, \texttt{model\_id}~$=0$), the racing
diffusion model (RDM, \texttt{model\_id}~$=1$), and the circular diffusion model
(CDM, \texttt{model\_id}~$=2$). Each family is embedded through a learned model
embedding table and shares a common encoder--decoder backbone; posterior draws
are produced by a conditional flow-matching decoder.

The experiment documented here asks a targeted transfer question:

\begin{quote}
\emph{Given a CogFormer already trained on $\{\text{DDM},\text{RDM},\text{CDM}\}$,
can we add a fourth family --- the Linear Ballistic Accumulator (LBA,
\texttt{model\_id}~$=3$) --- by fine-tuning, without retraining from scratch and
without degrading inference on the original three families?}
\end{quote}

The interest is twofold. Scientifically, it tests whether the amortization
backbone has learned a family-general representation of the shared
generalized-linear design space, such that a new family is reachable by a small
adaptation. Practically, it establishes a cheap protocol for growing the model
class one family at a time.

%============================================================================
\section{Why LBA is a favorable first target}
%============================================================================

LBA is chosen as the first fine-tuning target because it is maximally aligned
with a family the base model already knows, RDM. Concretely:

\begin{itemize}
  \item \textbf{Identical parameter space.} LBA and RDM share the same six
    intrinsic parameters under the shared response-coding convention: a base
    drift $v$, a drift advantage $v_{\mathrm{diff}}$ for the favored
    accumulator, a start-point upper bound $a$, non-decision time $\tau$, and
    the dispersions $s_v, s_\tau$. Both also carry the global contamination
    parameter $\pi$.
  \item \textbf{Identical intrinsic priors.} The prior over these parameters is
    literally the same function for the two families (\texttt{lba\_priors}
    reproduces \texttt{rdm\_priors} term for term). The two families therefore
    occupy the same region of parameter space and differ only in the generative
    mechanism that maps parameters to response times and choices.
  \item \textbf{Different likelihood geometry.} What differs is the accumulation
    process. LBA is \emph{ballistic}: within a trial there is no moment-to-moment
    noise; a linear accumulator with a random start point (uniform on $[0,a]$)
    and a Gaussian-perturbed but trial-constant drift races to a threshold
    $a + b$ with the class-level offset $b=0.5$. This yields a response-time
    distribution of a different shape from RDM's diffusive racers, even though
    the parameters carry the same names.
\end{itemize}

The upshot is that the adaptation required of the backbone is expected to be
\emph{small in parameter space but non-trivial in likelihood shape} --- an ideal
regime for fine-tuning, and the reason LBA's embedding is warm-started from
RDM's (Section~\ref{sec:setup}).

%============================================================================
\section{Setup and initialization}\label{sec:setup}
%============================================================================

\paragraph{Global parameter space.}
The base run amortizes all families in a shared \emph{global} parameter space:
each family's local coefficient matrix $\rmB_\clF$ is lifted, via
\texttt{lift\_to\_global\_space}, into a fixed-width global layout whose columns
are the union of all families' intrinsic parameters (with the global nuisance
$\pi$ appended). Fine-tuning preserves this layout exactly. All four families
--- including LBA --- are instantiated when constructing the \texttt{ModelClass}
so that the global column ordering (insertion order
DDM$\to$RDM$\to$CDM$\to$LBA) and hence the global parameter count matches the
base run bit-for-bit.

\paragraph{Loading the checkpoint.}
The pretrained checkpoint is loaded with \texttt{strict=True}; every backbone
weight transfers unchanged. The single structural change is the model embedding
table, which grows from three rows to four:
\[
  \text{model\_embedding} \in \mathbb{R}^{3 \times d}
  \;\longrightarrow\;
  \mathbb{R}^{4 \times d}, \qquad d = \texttt{model\_embed\_dim}.
\]
The new row 3 (LBA) is \textbf{warm-started from row 1 (RDM)} rather than
randomly initialized, exploiting the shared parameter space and priors so that
fine-tuning begins from a sensible family representation instead of noise.

\paragraph{Free and fixed intrinsics.}
During fine-tuning the free LBA intrinsics are $\{v, v_{\mathrm{diff}}, a,
\tau\}$; the dispersions $\{s_v, s_\tau\}$ are held at their fixed defaults
($0$ on the link scale). The architecture hyperparameters (encoder/decoder
depth and width, seed set, embedding dimensions, and crucially the
\texttt{self\_attn\_padding\_mask} flag) must match the base checkpoint exactly;
a mismatch on the padding-mask regime forces the network to re-adapt and
squanders the transfer.

%============================================================================
\section{Fine-tuning schemes}
%============================================================================

Two schemes are implemented, differing only in the data mixture presented during
fine-tuning.

\paragraph{Scheme 1 --- LBA-only.}
Each step samples only LBA data, lifts it to the global space, and tags it with
\texttt{model\_id}~$=3$. This is the fastest route to LBA competence but exposes
the shared backbone to LBA gradients exclusively, risking \emph{catastrophic
forgetting} of DDM/RDM/CDM.

\paragraph{Scheme 2 --- joint replay (default).}
Each step samples jointly from all four families through the
\texttt{ModelClass}, so DDM/RDM/CDM examples are continually replayed alongside
LBA. This preserves the original families by construction, at the cost of
slower LBA-specific progress per unit compute.

\paragraph{Optimization.}
Fine-tuning uses AdamW with a small learning rate ($1\times10^{-5}$, an order of
magnitude below pretraining), a cosine-annealing schedule over the epoch budget,
gradient-norm clipping at $5.0$, and light dropout ($0.05$). Nothing is frozen:
the full backbone remains trainable in both schemes.

%============================================================================
\section{Considered alternative: parameter-efficient adaptation}
%============================================================================

We considered low-rank adaptation (LoRA) and its quantized variant (QLoRA) as
alternative fine-tuning strategies. The conclusion is that \textbf{QLoRA is not
relevant} here and \textbf{LoRA is relevant only after reframing}.

\paragraph{QLoRA is inapplicable.}
QLoRA's sole purpose is to fit a base model too large to hold in full precision
by storing its frozen weights in low precision. CogFormer is small (a few
million parameters); there is no memory wall, so quantization offers no benefit
and only costs precision.

\paragraph{LoRA, reframed as regularization rather than efficiency.}
LoRA bundles two ideas: (i) freeze the base weights and learn only an additive
delta, and (ii) constrain that delta to be low rank, $\Delta\rmW = \mathbf{B}
\mathbf{A}$. At LBA's scale the efficiency motivation is void --- full
fine-tuning is already cheap --- so the low-rank constraint stops being a
cost-saver and becomes an \emph{inductive bias}: the hypothesis that the LBA
adaptation lives in a low-dimensional subspace of weight space. Given the shared
RDM parameter space, this hypothesis is plausible and would, if it holds, buy
two things: regularization against overfitting on limited LBA coverage, and a
\emph{toggleable} forgetting guarantee --- with the adapter disabled the network
reproduces the pretrained checkpoint exactly, so DDM/RDM/CDM inference is
preserved without replay.

\paragraph{Caveat.}
The low-rank delta acts on dense layers shared by every family's forward pass.
The non-forgetting guarantee therefore holds only if the other families are run
with the adapter disabled, or if the adapter is conditioned on
\texttt{model\_id} (a conditional-LoRA design). This is genuine implementation
work, not the one-line API of large-language-model tooling. A frozen-backbone,
low-rank ``Scheme 3'' is thus recorded as a promising comparison against the
current full-fine-tuning schemes, with its payoff riding entirely on how close
LBA sits to the RDM subspace.

%============================================================================
\section{Validation and forgetting audit}
%============================================================================

A dedicated validation script evaluates a fine-tuned checkpoint on two axes:

\begin{enumerate}
  \item \textbf{LBA competence.} Posterior recovery for LBA under its new
    embedding (\texttt{model\_id}~$=3$), across the benchmark design conditions
    (intercept-only, fixed, regressed, interaction, mixed-predictor, and
    held-out multi-factor designs).
  \item \textbf{Forgetting audit.} Recovery for DDM, RDM, and CDM under their
    \emph{original} embeddings ($0/1/2$), quantifying any degradation relative
    to the base checkpoint. Scheme~1 is expected to show measurable forgetting
    here; Scheme~2 is expected to hold.
\end{enumerate}

For each family $\times$ design condition the script emits saved posterior draws
(\texttt{*\_pred.npz}), a metrics table (\texttt{*\_metrics.csv}), and the
standard diagnostic panels: parameter recovery, simulation-based calibration
coverage, ECDF calibration, and per-parameter metric summaries. Global
parameters (e.g.\ $\pi$) are estimated jointly with the coefficient matrix but
plotted and scored separately.

%============================================================================
\section{Summary}
%============================================================================

The experiment adds LBA to a trained three-family CogFormer by (i)~extending the
model embedding table by one RDM-warm-started row, (ii)~loading all other weights
unchanged, and (iii)~fine-tuning either on LBA alone (Scheme~1) or on the full
four-family mixture (Scheme~2, default). LBA is a deliberately favorable first
target: it shares RDM's parameter space and priors exactly while differing in
likelihood geometry, so the required adaptation is small in parameter space but
real in distributional shape. Success is measured jointly by LBA recovery and by
the absence of forgetting on the original three families. A frozen-backbone
low-rank adapter (LoRA-style, conditioned on \texttt{model\_id}) is identified as
the natural next comparison; quantization (QLoRA) is ruled out as irrelevant at
this scale.

\end{document}
